% Generado por R/exportacion.R desde reports/Entrega 3.qmd
\begingroup\tiny\setstretch{1}
\par
\noindent\begin{minipage}{\textwidth}
\centering
\captionof{table}{Varianza del puntaje global asociada a cada variable: de forma aislada, aporte propio dentro del modelo conjunto y parte compartida con las demás.}\label{tab:aportes}
\begin{tabular}{l r r r}
  \toprule
  \textbf{Variable} & \textbf{Aislada} & \textbf{Propia} & \textbf{Compartida} \\
  \midrule

  \multicolumn{4}{l}{\textit{Características individuales}} \\
  Edad & 9,43 & 2,23 & 7,20 \\
  Horas de trabajo & 4,79 & 1,04 & 3,75 \\
  Dedicación a la lectura & 2,51 & 1,39 & 1,12 \\
  Género & 0,52 & 0,91 & $-$0,39 \\
  \multicolumn{4}{l}{\textit{Contexto del hogar}} \\
  Educación de la madre & 12,63 & 1,66 & 10,97 \\
  Libros en el hogar & 9,32 & 1,18 & 8,13 \\
  Acceso a internet & 9,23 & 0,58 & 8,66 \\
  Estrato de la vivienda & 7,19 & 0,37 & 6,82 \\
  \multicolumn{4}{l}{\textit{Establecimiento educativo}} \\
  Jornada del colegio & 9,84 & 1,71 & 8,14 \\
  Naturaleza del colegio & 6,43 & 0,40 & 6,03 \\
  Zona del colegio & 3,75 & 0,54 & 3,21 \\
  \multicolumn{4}{l}{\textit{Geografía}} \\
  Región & 4,60 & 1,45 & 3,16 \\
  \bottomrule
\end{tabular}
\par\raggedright{\footnotesize Porcentajes de varianza. El modelo conjunto explica el 33,0\,\%, la suma de las asociaciones aisladas alcanza el 80,2\,\% y la diferencia es redundancia entre variables. Una parte compartida negativa indica supresión.\par}
\end{minipage}
\par
\endgroup
